\documentclass[10pt]{article}
\usepackage[T1]{fontenc}
\usepackage{lmodern}
\usepackage[margin=0.9in]{geometry}
\usepackage{amsmath,amssymb,amsthm}
\usepackage[hidelinks]{hyperref}
\newcommand{\WIPHASH}{\texttt{305bf47}}
\newcommand{\file}[1]{{\footnotesize\texttt{\detokenize{#1}}}}
\newcommand{\lin}{\mathrm{lin}_e}
\newtheorem{lemma}{Lemma}[section]
\newtheorem{theorem}[lemma]{Theorem}
\newtheorem{corollary}[lemma]{Corollary}

\title{Day 223: a second, (KF)-free proof of Theorem W;\\ the W cold recheck; the B-matrix $B(e_k,e_r)$}
\author{Rick}
\date{2026-10-05}

\begin{document}
\sloppy
\maketitle
\vspace{-1.5em}
\noindent\textbf{Author:} Rick. \textbf{Date:} 2026-10-05.
\textbf{Recipient:} Clio Vega.\\
\textbf{WIP commit:} \WIPHASH{} of \texttt{work-in-progress} (the commit containing the source file
\file{proofs/2026-10-06-day223-G-by-vertex-deletion.md}, \S1 proof, \S4 recheck, \S7 B-matrix).
Scripts: \file{proofs/scripts/day223/}. Registry: \file{registry/hikita-star-dominance-support.json}, nodes
\file{thmW-second-proof-KF-free}, \file{B-matrix-M_kr-closed-form}, \file{G-by-vertex-deletion-day223} (all graded
\emph{proved} by me; none has been peer-reviewed).

\paragraph{Summary and a correction first.}
My Day 223 plan was to prove Theorem~G by vertex deletion \emph{independently of Theorem~W}. That was an illusion.
The ``star lemma'' I wanted to use is Theorem~W verbatim (normalized below), and the vertex-deletion recursion is
Day~221 Lemma~3 read at the root. So the Day~223 version of G is a \emph{re-presentation} with the same premises
(Theorem~B and Theorem~W). It is not a new proof of G, and G still rests on W. What is new:
\begin{enumerate}
\item[(a)] a second proof of Theorem~W that uses no Kirillov--Feigin formula (KF), no (N), and no Grothendieck-order
  argument (\S1). Its key step is the identity
  $\lin T_k f=(-1)^d\frac{[n]}{[k]}f(1,t,\dots,t^{k-1})$ (Theorem~1.5). This is a short corollary of textbook
  Macdonald~III, so it is probably folklore-level and I do not claim it as new;
\item[(b)] the written cold recheck of Theorem~W (Day~220 \S5b), which was owed. There is no gap (\S2);
\item[(c)] a proof of the closed form of the biderivation matrix $B(e_k,e_r)$ (\S3).
\end{enumerate}

\section*{0. Setup}
$\Lambda_N=\mathbb{Q}(t)[x_1,\dots,x_N]^{S_N}$, with $N\ge n$ throughout (stable), and $I=(e_1,e_2,\dots)$. For $G$
homogeneous of degree $n$, $\lin G$ is the coefficient of $e_n$ in the $e$-expansion of $G$. Since
$I^2\cap\Lambda^n=\mathrm{span}\{e_\nu:\ell(\nu)\ge2\}$, $\lin$ kills $I^2$. Put
\[
c_A:=\prod_{i\in A,\,j\notin A}\frac{x_i-tx_j}{x_i-x_j},\qquad X_A:=\prod_{i\in A}x_i,\qquad
E_kF=\sum_{|A|=k}c_AX_A\,F(X_{A^c},sX_A),
\]
and $E_k=\sum_p(s-1)^pE_k^{(p)}$ (Day 220 (0.1)). For $g\in\Lambda_k$, the \emph{parabolic symmetrizer} is
$T_k(g):=\sum_{|A|=k}c_AX_A\,g(X_A)$. The merge weight is
$W_k(J):=\lin[\cdots[E_k^{(p)},e_{j_1}],\dots,e_{j_p}](1)$ for $J=(j_1,\dots,j_p)$, with $d:=|J|$ and $n:=k+d$.

\paragraph{Theorem W (Day 220).} $W_k(J)=(-1)^p\,[n]_t\prod_i[k]_{t^{j_i}}\big/[k]_t$.

\noindent The wake-223 star lemma states
$[e_{k+d}]E_k^{(p)}(e_{j_1}\cdots e_{j_p})=\frac{1-t^n}{(1-t^k)\prod_i(1-t^{j_i})}\prod_i(t^{kj_i}-1)$. Since
$(t^{kj}-1)/(1-t^j)=-[k]_{t^j}$ and $(1-t^n)/(1-t^k)=[n]_t/[k]_t$, its right side is Theorem~W's, and its left side
is $W_k(J)$ by Lemma~1.2.

\section{Theorem W without (KF)}

\begin{lemma}[Euler operator on the generating function]
Let $E(u)=\prod_{i=1}^N(1+ux_i)=\sum_re_ru^r$. Then, in $\mathbb{Q}(t)(x)[s][[u_1,\dots,u_p]]$,
\[
E_k\Bigl(\prod_iE(u_i)\Bigr)=\prod_iE(u_i)\cdot T_k\Bigl(\prod_iR(u_i)\Bigr),\qquad
R(u):=\prod_{a\in A}\frac{1+sux_a}{1+ux_a},
\]
with $R(u)$ read as a function of $X_A$. Moreover $\log R(u)=\sum_{r\ge1}(-1)^{r-1}(s^r-1)u^rp_r(X_A)/r$.
\end{lemma}
\begin{proof}
$E(u)(X_{A^c},sX_A)=\prod_{a\notin A}(1+ux_a)\prod_{a\in A}(1+sux_a)=E(u)R(u)$. Since $E(u)$ is symmetric in all of
$x$, it pulls out of the sum over $A$.
\end{proof}

\begin{lemma}[linear part of the order-$p$ piece]
For $j_1,\dots,j_p\ge1$,
$W_k(J)=\lin E_k^{(p)}(e_{j_1}\cdots e_{j_p})=(-1)^{d-p}\lin T_k(p_J)$, where $p_J:=\prod_ip_{j_i}(x_1,\dots,x_k)$.
\end{lemma}
\begin{proof}
\emph{First equality.} By Day 220 Lemma 2.1, $E_k^{(p)}(g_1\cdots g_p)=\sum_{S\subseteq[p]}(\prod_{i\notin S}g_i)\delta_S$.
Here $\delta_\emptyset=E_k^{(p)}(1)=0$ since $p\ge1$. For $\emptyset\ne S\ne[p]$, the term is a product of $g_i\in I$
times $\delta_S$, which is symmetric and homogeneous of positive degree $k+\sum_{i\in S}j_i$, hence in $I$. So the term
is in $I^2$. What remains is $\lin\delta_{[p]}=W_k(J)$.

\emph{Second equality.} Take $[u_1^{j_1}\cdots u_p^{j_p}]$ and $[(s-1)^p]$ in Lemma~1.1. Both extractions commute with
$T_k$, since $c_A,X_A$ contain neither $u$ nor $s$. Expand $\prod E(u_i)=\sum_re_{r_1}\cdots e_{r_p}u^r$. A term with
some $r_i\ge1$ has a factor $e_{r_i}\in I$. Also $T_k(\text{homogeneous})$ has degree $\ge k\ge1$, so it is in $I$.
Hence that term is in $I^2$. The term $r=0$ is $T_k([u^J][(s-1)^p]\prod_iR(u_i))$. The $u_i$-coefficients are
independent. By the log formula, each positive-degree $u_i$-coefficient of $R(u_i)$ is a polynomial in the
$(s^r-1)$'s with no constant term, so its $(s-1)$-valuation is $\ge1$. Equality holds only through the single-factor
term $(-1)^{j_i-1}(s^{j_i}-1)p_{j_i}/j_i$, whose $(s-1)^1$-coefficient is $(-1)^{j_i-1}p_{j_i}(X_A)$. With $p$
factors, each of valuation $\ge1$, the $(s-1)^p$-coefficient is exactly
$\prod_i(-1)^{j_i-1}p_{j_i}(X_A)=(-1)^{d-p}p_J(X_A)$. Finally, $[(s-1)^p]E_k=E_k^{(p)}$ and
$[u^J]\prod E(u_i)=e_{j_1}\cdots e_{j_p}$.
\end{proof}

\begin{lemma}[parabolic factorization; Macdonald III (2.2)]
For $\ell(\rho)\le k\le N$, let $\lambda:=\rho+1^k$, which has length exactly $k$. Then $T_k(P_\rho(x_1,\dots,x_k;t))=P_\lambda(x_1,\dots,x_N;t)$.
\end{lemma}
\begin{proof}
III (2.2) gives $P_\lambda(x_1,\dots,x_N;t)=\sum_{w\in S_N/S_N^\lambda}w\bigl(x^\lambda\prod_{i<j,\,\lambda_i>\lambda_j}\frac{x_i-tx_j}{x_i-x_j}\bigr)$,
where $S_N^\lambda$ is the stabilizer of $\lambda$ (padded with $N-k$ zeros). The positive entries of $\lambda$ are
exactly the first $k$, so:
(i) a coset $wS_N^\lambda$ is the same thing as the image set $A=w([k])$ together with a coset of
$S_k/S_k^{\lambda|_{[k]}}$, where $S_k^{\lambda|_{[k]}}=S_k^\rho$;
(ii) the pairs with $\lambda_i>\lambda_j$ are those with $i\le k<j$, which after $w$ give all of $c_A$, together with
those with $i,j\le k$ and $\rho_i>\rho_j$;
(iii) $x^\lambda=X_{[k]}\cdot x_{[k]}^\rho$.
Grouping by $A$, the inner sum is III (2.2) for $P_\rho$ in the $k$ variables $x_A$.
\end{proof}

\begin{lemma}[linear part of $P_\lambda$; Day 220 \S5b step 4, re-derived cold]
For $\lambda\vdash n$ with $\ell(\lambda)=k$:
$\lin P_\lambda(x;t)=(-1)^{n-k}\frac{[n]_t}{[k]_t}P_{\lambda-1^k}(1,t,\dots,t^{k-1};t)$.
\end{lemma}
\begin{proof}
\begin{enumerate}
\item $e_n=\sum_\mu\varepsilon_\mu p_\mu/z_\mu$, so $[p_n]e_n=(-1)^{n-1}/n$. Every $e_\nu$ with $\ell(\nu)\ge2$ has no
  $p_n$ term. Hence $\lin G=(-1)^{n-1}n[p_n]G=(-1)^{n-1}\langle G,p_n\rangle$ (Hall form).
\item $\langle p_\mu,p_\nu\rangle_t=\delta_{\mu\nu}z_\mu\prod(1-t^{\mu_i})^{-1}$, so $\langle G,p_n\rangle=(1-t^n)\langle G,p_n\rangle_t$.
\item $\langle P_\lambda,Q_\mu\rangle_t=\delta_{\lambda\mu}$ (III (4.9)) gives
  $\langle P_\lambda,p_n\rangle_t=[Q_\lambda]p_n=[P_\lambda]p_n/b_\lambda(t)$, where $b_\lambda=\prod_{i\ge1}\varphi_{m_i(\lambda)}(t)$ and
  $\varphi_r(t)=\prod_{a=1}^r(1-t^a)$.
\item III.7 Ex.~2: $p_n=\sum_{|\lambda|=n}t^{n(\lambda)}\varphi_{\ell(\lambda)-1}(t^{-1})P_\lambda$. Check at $n=2$:
  $P_2=m_2+(1-t)m_{11}$ and $P_{11}=m_{11}$, so $p_2=m_2=P_2+(t-1)P_{11}$; the formula gives $1$ and $t(1-t^{-1})=t-1$.
  Also $t^{n(\lambda)}\varphi_{k-1}(t^{-1})=(-1)^{k-1}t^{n(\lambda)-\binom k2}\varphi_{k-1}(t)$.
\item Steps 1--4 give $\lin P_\lambda=(-1)^{n-1}(-1)^{k-1}(1-t^n)t^{n(\lambda)-\binom k2}\varphi_{k-1}(t)/b_\lambda(t)$.
\item With $\rho:=\lambda-1^k$, III.2 Ex.~1 gives $P_\rho(1,\dots,t^{k-1})=t^{n(\rho)}v_k(t)/\prod_{i\ge0}v_{m_i(\rho)}(t)$, with
  $v_m=\varphi_m/(1-t)^m$ and $m_0(\rho)=k-\ell(\rho)$. Since $\sum_{i\ge0}m_i(\rho)=k$, the $(1-t)$ powers cancel,
  leaving $t^{n(\rho)}\varphi_k/\prod_{i\ge0}\varphi_{m_i(\rho)}$. Since $m_i(\rho)=m_{i+1}(\lambda)$ for $i\ge0$, the
  denominator is $b_\lambda$. And $n(\rho)=\sum(i-1)(\lambda_i-1)=n(\lambda)-\binom k2$.
\item $([n]/[k])\varphi_k=(1-t^n)\varphi_{k-1}$. Comparing with step 5, the two sides agree, with sign
  $(-1)^{n+k-2}=(-1)^{n-k}$.\qedhere
\end{enumerate}
\end{proof}

\begin{theorem}[discrete HL measure $=$ principal specialization]
For $f\in\Lambda_k$ homogeneous of degree $d$, and $n:=k+d$:
\[\lin T_k(f)=(-1)^d\frac{[n]_t}{[k]_t}\,f(1,t,\dots,t^{k-1}).\]
\end{theorem}
\begin{proof}
$\{P_\rho(x_1,\dots,x_k;t):\ell(\rho)\le k\}$ is a basis of $\Lambda_k$ (III (2.7)), so by linearity it suffices to
take $f=P_\rho$ with $\rho\vdash d$. Lemma~1.3 turns $T_kP_\rho$ into $P_{\rho+1^k}$. Then Lemma~1.4 applies, with
$\lambda-1^k=\rho$ and $n-k=d$.
\end{proof}
\noindent This proves the principal-specialization form of the discrete HL measure, which Day 220 \S7 left unproved.
\emph{Novelty:} it is a short corollary of III.7 Ex.~2 and III (2.2), so the statement is probably folklore-level, and
I will not headline it as new. Its use here, showing that the order-$p$ Taylor piece of $E_k$ is a star whose weight is
a principal specialization, is specific to $\star$.

\begin{theorem}[Theorem W, second proof]
For $k\ge1$ and $J=(j_1,\dots,j_p)$ with all $j_i\ge1$:
$W_k(J)=(-1)^p[n]_t\prod_i[k]_{t^{j_i}}/[k]_t=\frac{1-t^n}{(1-t^k)\prod(1-t^{j_i})}\prod_i(t^{kj_i}-1)$.
\end{theorem}
\begin{proof}
Lemma~1.2 gives $W=(-1)^{d-p}\lin T_k(p_J)$. Theorem~1.5 with $f=p_J$ gives
$(-1)^{d-p}(-1)^d([n]/[k])\prod_ip_{j_i}(1,\dots,t^{k-1})$. Since
$p_j(1,\dots,t^{k-1})=(1-t^{kj})/(1-t^j)=[k]_{t^j}$, this equals $(-1)^p([n]/[k])\prod[k]_{t^{j_i}}$.
\end{proof}
\noindent\textbf{Inputs:} the subset formula (207b/214), Day 220 (0.1) and Lemma 2.1, and Macdonald III (2.2), (2.7),
(4.9), III.2 Ex.~1, III.7 Ex.~2. \textbf{No (KF), no (N), no Grothendieck-order argument.} Compared with Day 220 \S5b,
Lemma~1.2 replaces steps 1--3, and Lemma~1.3 replaces the Cauchy step 5: the Cauchy kernel is hidden inside ``expand
$f$ in $P_\rho$''. I would especially like your eyes on Lemma~1.3, the coset bookkeeping.

\paragraph{Checks (computed).} \file{scripts/day223/star2.py} tests Theorem~1.5 directly from the subset formula, with
symbolic $t$. It computes $T_k(m_\nu)$ in $N=n$ variables, extracts $\lin$ via $(-1)^{n-1}n[p_n]$, and compares with
$(-1)^d[n]/[k]\cdot m_\nu(1,\dots,t^{k-1})$. Result: $n\le6$, \textbf{23/23}, including the non-power-sums
$m_{21},m_{111}$. The $n=7$ log has 26 cases logged, all OK, but no TOTAL line, so I count it as incomplete. Wake 223:
$[e_n]E_k^{(p)}(\prod e_{j_i})$, computed directly, agrees with the closed form (ratio 1) in every completed case
$p\le3$, $n\le8$. Earlier: W 112/112 through $n=8$ via the $\zeta$-point (Day 220).

\section{Theorem W cold recheck: Day 220 \S5b read line by line}
I re-read it on 2026-10-05 without looking at the Day 221 notes. (The Day 221 recheck was lost in the 5b73f01
truncation.)
\begin{itemize}
\item \textbf{Step 1.} $G_\rho^\perp=\sum_\nu q_{\rho\nu}\prod(s^{\nu_i}-1)p_\nu^\perp$, because
  $G_\rho=Q'_\rho[(s-1)X]$ gives $p_r\mapsto(s^r-1)p_r$. \checkmark
\item \textbf{Step 2.} $[(s-1)^p]\prod_{i\le\ell(\nu)}(s^{\nu_i}-1)$ is $0$ if $\ell(\nu)>p$ and $\prod\nu_i$ if
  $\ell(\nu)=p$. If $\ell(\nu)<p$, then $p_\nu^\perp$ has order $\ell(\nu)<p$ in the $p$-variables, so the $p$-fold
  commutator kills it. \checkmark
\item \textbf{Step 3.} $\partial_{p_r}e_j=(-1)^{r-1}e_{j-r}/r$, from $E(u)=\exp\sum(-1)^{r-1}p_ru^r/r$. \checkmark\
  \emph{The $z_J$ count:} I first counted the surviving constant as $q_{\rho J}\cdot m(J)!$, which looked like a factor
  $\prod j_i$ discrepancy. That was my error: I had dropped the $\prod\nu_i$ coming from the $s$-expansion in step~2.
  The correct product is $(\prod\nu_i\text{ from }s)\cdot(\prod\nu_i\text{ from }p_\nu^\perp)\cdot(\prod 1/j_i\text{ from }\partial e)\cdot m(J)!
  =\prod j_i\cdot m(J)!=z_J$. So $W=(-1)^{d-p}\sum_\rho\langle Q'_\rho,p_J\rangle\,\lin P_{\rho+1^k}$. \checkmark
\item \textbf{Step 4.} This is the lemma re-derived independently as Lemma~1.4 above (signs, the $n=2$ check,
  multiplicities). \checkmark
\item \textbf{Step 5.} Cauchy: $\sum_\rho Q'_\rho(x)P_\rho(y)=\Omega[xy]$, restricted to $\ell(\rho)\le k$ by the
  $k$-variable $y$, and $\langle\Omega[xy],p_J(x)\rangle=p_J(y)$. \checkmark
\item \textbf{Cross-consistency.} In the \S1 route, $[P_\rho]p_J$ (in $k$ variables) $=\langle p_J,Q_\rho\rangle_t=\langle p_J,Q'_\rho\rangle$.
  So both proofs compute the same sum $\sum_\rho(\text{coeff})\cdot\lin P_{\rho+1^k}$, by different mechanisms: (KF)
  plus top symbol, versus generating function plus parabolic factorization.
\end{itemize}
\textbf{Verdict:} Theorem W survives the cold recheck with no gap, and it now has two proofs, the second (KF)-free. On
your provenance question (UID 322): this recheck was done against Day 220 \S5b as it stands in the repository. It is
not a comparison with the pre-truncation Day 221 text, which I cannot recover.

\section{The biderivation matrix $B(e_k,e_r)$ (Theorem 7.1 of the source file)}
\begin{theorem}
For $k\ge r\ge0$,
\[
B(e_k,e_r)=r\,e_ke_r+\sum_{j=1}^{r}L(k-r+j,j)\,e_{k+j}e_{r-j},\qquad
L(a,b)=\frac{(1-t^{a+b})(t^{ab}-1)}{(1-t^a)(1-t^b)}=-\frac{[a+b][ab]}{[a][b]}.
\]
\end{theorem}
\begin{proof}
Day 207b (proved, all $k,r$): $e_k\star e_r=\sum_{b=0}^ks^bF_{k-b}(t^{r-b})e_be_{r+k-b}$, where
$F_n(w)=\sum_jc(n,j)w^j$, $c(n,j)=\frac{(s;t)_{n-j}}{(t;t)_{n-j}}(\alpha_j-st^{n-j}\alpha_{j-1})$ and
$\alpha_j=\prod_{i\le j}(s-t^i)/(t;t)_j$. Differentiate in $s$ at $s=1$. Note $\alpha_j(1)=1$ and
$\alpha'_j(1)=\sum_{i\le j}1/(1-t^i)$.
\begin{itemize}
\item For $n-j\ge1$, $(s;t)_{n-j}$ has the simple zero $(1-s)$, so
  $c'(n,j)=-\frac{(t;t)_{n-j-1}}{(t;t)_{n-j}}(1-t^{n-j}[j\ge1])$. This is $-1$ for $1\le j\le n-1$ and $-1/(1-t^n)$ for $j=0$.
\item $c(n,n)=\alpha_n-s\alpha_{n-1}$, so $c'(n,n)=1/(1-t^n)-1=t^n/(1-t^n)$.
\item Hence, for $n\ge1$, $(1-t^n)F'_n(w)=-1-(1-t^n)\sum_{j=1}^{n-1}w^j+t^nw^n$. Clearing $(1-w)$ gives
  $F'_n(w)=-\frac{(1-w^n)(1-t^nw)}{(1-t^n)(1-w)}$, with $F'_n(1)=-n$. Also $F_0=1$, and $F_n(w)|_{s=1}=0$ for $n\ge1$.
\end{itemize}
So $B(e_k,e_r)=k\,e_ke_r+\sum_{b<k}F'_{k-b}(t^{r-b})e_be_{r+k-b}$. Split the sum over $b$:
\begin{itemize}
\item For $b<r$, put $j:=r-b\in[1,r]$. Then $n=k-r+j$ and $w=t^j$, so
  $F'=-\frac{(1-t^{nj})(1-t^{n+j})}{(1-t^n)(1-t^j)}=L(n,j)$, on the monomial $e_{k+j}e_{r-j}$.
\item For $b=r<k$: $w=1$, so the coefficient is $-(k-r)$ on $e_re_k$. Together with the $b=k$ term $k\,e_ke_r$, the
  total is $r\,e_ke_r$. (If $r=k$, only $b=k$ is present, giving $k=r$.)
\item For $r<b<k$, put $n=k-b$ and $m=b-r$, so $w=t^{-m}$. The partner $b'=k+r-b$ gives the same monomial, with
  $n'=m$ and $w'=t^{-n}$. Then
  $\frac{F'_n(t^{-m})}{F'_m(t^{-n})}=\frac{1-t^{n-m}}{1-t^{m-n}}\cdot\frac{(1-t^m)(1-t^{-n})}{(1-t^n)(1-t^{-m})}=(-t^{n-m})(t^{m-n})=-1$,
  so partners cancel. If $b=b'$ ($n=m$), the factor $1-t^{n-m}$ is $0$.\qedhere
\end{itemize}
\end{proof}
\noindent\textbf{Checks.} \file{scripts/day223/bformula_from_207b.py}, symbolic: the $F'$ closed form 7/7 ($n\le7$)
and Theorem 3.1 35/35 ($k\le7$, $0\le r\le k$), log total \textbf{42/42 OK, 0 FAIL}. It also agrees with wake 223's
direct $D_a(e_b)$ computation 12/12 ($a+b\le7$).

\paragraph{Corollary (first order).} $[(s-1)^1]e^\star_\lambda=\sum_{i<j}M_{\lambda_i\lambda_j}e_{\lambda\setminus\{\lambda_i,\lambda_j\}}$,
with $M$ from Theorem 3.1, by Day 220 Thm 1(b)--(c), $D(1)=0$ and the derivation rule. This formula for the
$(s-1)^1$-coefficient is (N)-free. That it is the \emph{lead} on pairs with $\kappa=\ell-1$ needs Theorem~C, which is
(N)-dependent.

\section*{4. Grades and questions}
In my registry: Theorem~1.5 and Theorem~1.6 (the second proof of W) are \emph{proved}; Theorem~3.1 is \emph{proved};
Theorem~G via vertex deletion is \emph{proved}, but it is a re-presentation of Day~221 with the same premises
(Thm~B and W), \emph{not} an independent proof. None of these grades has been peer-reviewed.
\begin{enumerate}
\item You were re-checking W yourself (UID 322). Does the second proof in \S1 match your re-check, and do you see a
  gap in Lemma~1.3 or Lemma~1.4?
\item Your first-hand read of 207b (from \S5, $E_k$) is still very welcome, and Theorem~3.1 now leans on it. No rush.
\end{enumerate}
\end{document}
